\chapter{Typesetting software citations in venue classes}
\label{app:venue-classes}

These notes are set in the plain \texttt{report} class, where loading
\texttt{biblatex-software} takes two lines. Submitting to a conference or a journal is less
forgiving: the major venue classes ship their own bibliography machinery. The Hands-on box in
Chapter~\ref{ch:ardc} asks you to drop your software entries into a real venue template; this
appendix is the part you would otherwise have to reverse-engineer.

One rule holds everywhere: build with \texttt{biber}, never with \texttt{bibtex}. Beyond that there
are two situations, and telling them apart saves a great deal of confusion. Either the class knows
nothing about software entries --- in which case you load \texttt{biblatex} with
\texttt{datamodel=software} and then \texttt{software-biblatex} yourself, as for \texttt{IEEEtran}
below --- or the class \emph{already integrates} \texttt{biblatex-software}, in which case you must
\emph{not} load it again. \texttt{acmart} is now the second kind, which is why it needs less
configuration than folklore suggests, not more.

\section{ACM: \texttt{acmart}}

This used to be the awkward case, and much advice still found online says so: \texttt{acmart}
preloads \texttt{natbib}, which collides with \texttt{biblatex}, and for years the remedy was a
hand-rolled incantation with macros to unbind. That is no longer necessary. Since version~1.84 the
class ships a rewritten \texttt{biblatex} layer and two bibliography styles of its own ---
\texttt{acmnumeric} and \texttt{acmauthoryear} --- which reproduce the behaviour of the traditional
\texttt{ACM-Reference-Format.bst}. Better still for our purposes, those styles \emph{already
integrate} \texttt{biblatex-software}: \texttt{acmdatamodel.dbx} pulls in the software data model,
and \texttt{acmnumeric.bbx} pulls in the software bibliography driver and switches \texttt{swhid} on
for you.

The consequence is that you do \emph{less} here than in a plain document, not more. You do not pass
\texttt{datamodel=software}, and you do not load \texttt{software-biblatex} yourself --- the class
does both. What you must do is turn \texttt{natbib} off and pick an ACM style:

\begin{quote}\footnotesize
\begin{verbatim}
\documentclass[sigconf,natbib=false]{acmart}

\RequirePackage[
  datamodel=acmdatamodel,   % pulls in the biblatex-software data model
  style=acmnumeric,         % or acmauthoryear, if the venue requires it
  ]{biblatex}

\addbibresource{yourbib.bib}
\end{verbatim}
\end{quote}

\noindent Then put \texttt{\textbackslash printbibliography} where the references belong. If you are
converting an existing ACM paper, delete any \texttt{\textbackslash citestyle\{acmauthoryear\}} line
as well: it belongs to the \texttt{natbib} world you have just left.

\texttt{biblatex-software} must still be installed --- the class \emph{uses} it rather than
\emph{replacing} it --- but you no longer name it anywhere. All four entry types then render inside
the ACM reference list, each carrying its \textsc{swhid}: a project as \texttt{[SW]}, a release as
\texttt{[SW Rel.]}, a module as \texttt{[SW Mod.]}, and a code fragment as \texttt{[SW exc.]}, with
the \texttt{crossref} relations printed as \enquote{\emph{from}} and \enquote{\emph{part of}} just as
they are in these notes.

\begin{callout}{One caveat, from ACM itself}
\small The class documentation is blunt about the status of this path: the \texttt{biblatex} styles
are \enquote{highly experimental} and are \emph{not} supported by \textsc{taps}, the ACM production
system that builds the final camera-ready version. Read that as: excellent for preprints, for arXiv,
and for your own archival copy; but check with the venue before relying on it for camera-ready, and
be ready to fall back to \texttt{ACM-Reference-Format.bst} --- losing the software entry types in the
process --- if production insists.
\end{callout}

\section{IEEE: \texttt{IEEEtran}}

\texttt{IEEEtran} is the first kind of class: it knows nothing about software entries, but it also
does not preload a competing citation package, so \texttt{biblatex} loads in the ordinary way and you
supply the software support yourself.

\begin{quote}\footnotesize
\begin{verbatim}
\documentclass[conference]{IEEEtran}

\usepackage[backend=biber, style=ieee, datamodel=software]{biblatex}
\usepackage{software-biblatex}

\ExecuteBibliographyOptions{
  swhid=true, swlabels=true, vcs=true, license=true}

\addbibresource{yourbib.bib}
\end{verbatim}
\end{quote}

\noindent That is the whole of it. All four entry types render inside the IEEE numbered reference
list, each carrying its \textsc{swhid}, and the \texttt{crossref} chains behave as they do everywhere
else --- a code fragment prints as \enquote{\ldots{} \emph{from} [1]} under its parent release rather
than repeating the project's authors, repository, and license.

\section{Other classes}

For any class not covered here, first ask which of the two kinds it is. If it already loads
\texttt{natbib}, expect a collision and look for a \texttt{natbib=false} class option before reaching
for macro surgery. If it ships its own \texttt{biblatex} styles, check whether they already input
\texttt{software.dbx} and \texttt{software.bbx} --- if they do, load neither the software data model
nor the style yourself. And if the venue insists on \texttt{bibtex} and will not accept
\texttt{biber}, the software entry types are simply not available: fall back to \texttt{@misc} and
put the \textsc{swhid} in the \texttt{howpublished} field, which is lossy, but at least keeps a
verifiable identifier in the record.
